%-------------------------
% Resume in LaTeX
% Based on Jake's Resume style, adapted for Chinese CS/AI resume
%------------------------

\documentclass[letterpaper,12pt]{article}

\usepackage[UTF8]{ctex}
\usepackage[
  left=0.58in,
  right=0.58in,
  top=0.44in,
  bottom=0.44in
]{geometry}
\usepackage{titlesec}
\usepackage{enumitem}
\usepackage[hidelinks]{hyperref}
\usepackage{tabularx}
\usepackage{xcolor}

\setmainfont[
  Path = /usr/local/texlive/2025/texmf-dist/fonts/opentype/public/tex-gyre/,
  Extension = .otf,
  UprightFont = texgyreheros-regular,
  BoldFont = texgyreheros-bold,
  ItalicFont = texgyreheros-italic,
  BoldItalicFont = texgyreheros-bolditalic,
  Numbers = Lining,
  Ligatures=NoCommon
]{texgyreheros}

\ifdefined\pdfgentounicode
  \input{glyphtounicode}
  \pdfgentounicode=1
\fi

\pagestyle{empty}
\urlstyle{same}
\raggedright
\raggedbottom
\setlength{\tabcolsep}{0pt}
\setlength{\parindent}{0pt}
\linespread{1.05}

\titleformat{\section}
  {\vspace{4pt}\large\bfseries}
  {}
  {0pt}
  {}
  [\vspace{1pt}\color{black}\titlerule\vspace{-2pt}]

\newcommand{\entry}[4]{
  \begin{tabular*}{\textwidth}{l@{\extracolsep{\fill}}r}
    \textbf{#1} & #2 \\
    #3 & #4 \\
  \end{tabular*}\vspace{-3pt}
}

\newcommand{\project}[1]{
  \vspace{3pt}\textbf{#1}\vspace{-3pt}
}

\setlist[itemize]{
  leftmargin=1.35em,
  itemsep=2.6pt,
  topsep=2.6pt,
  parsep=0pt,
  partopsep=0pt
}

\newenvironment{projectitems}{
  \begin{itemize}[
    leftmargin=1.35em,
    itemsep=3.2pt,
    topsep=3pt,
    parsep=1pt,
    partopsep=0pt
  ]
}{
  \end{itemize}
}

\begin{document}
\fontsize{13.5pt}{16.7pt}\selectfont

%----------HEADING----------
\begin{center}
  {\Huge\bfseries 乔启凡} \\ \vspace{4pt}
  \normalsize
  (+86)19935166136 \quad | \quad
  \href{mailto:qiaoqifan@zju.edu.cn}{qiaoqifan@zju.edu.cn} \quad | \quad
  \href{https://github.com/Hugaqq}{github.com/Hugaqq}
\end{center}

%-----------EDUCATION-----------
\section{教育经历}
\entry
  {浙江大学 竺可桢学院 图灵班（计算机科学与技术专业）}
  {2025.9 -}
  {GPA: 4.48/5.00 \textbar{} 排名: 13/68}
  {}
\begin{itemize}
  \item 参与图灵学术班项目：从图灵班 68 人中选 4 人参加科研训练（2026.3 -）
\end{itemize}

%-----------RESEARCH-----------
\section{科研经历}
\entry
  {庄博涵老师 ZIP Lab}
  {2026.4 -}
  {}
  {}

%-----------PROJECTS-----------
\section{项目}
\project{Efficient Diffusion RL infra for Visual Generation}
\begin{projectitems}
  \item 为 Diffusion RL 视觉生成算法搭建统一训练框架，支持基于 Flow-GRPO 的 Diffusion RL 算法训练，支持不同的采样方式、Reward 算法和模型
  \item 将训练流程抽象为 Rollout、Manifest、Reward、Optimizer、Logger 五个模块，简化新算法训练 infra 的构建，降低新模型、新 reward 和新算法的接入成本
  \item 目前正在整合 TempFlow-GRPO、Flash-GRPO、World-R1，完成了最小单元测试
\end{projectitems}

\project{RV64I RISC-V CPU 设计}
\begin{projectitems}
  \item 基于 SystemVerilog 设计并实现了 RV64I 单周期处理器核心
  \item 完成 Datapath、Controller 以及各个数据通路部件（如 ALU、RegFile 等）的编写，使其可以支持 RISC-V 的非扩展类型指令
  \item 实现基于 valid-ready 握手协议的多周期控制 FSM，将取指、执行、访存和提交解耦
\end{projectitems}

\project{对于字符表达式的自动微分系统}
\begin{projectitems}
  \item 基于 C 语言设计并实现了一个符号自动求导程序，支持中缀代数表达式的解析与计算
  \item 完成词法分析、递归下降解析器、表达式树构建、符号求导和表达式化简等核心模块
  \item 支持多变量求偏导及常见数学函数求导规则
\end{projectitems}

%-----------SKILLS-----------
\section{技术栈}
\begin{itemize}
  \item \textbf{编程语言/框架：}Python, PyTorch, C, Bash, Verilog/SystemVerilog
  \item \textbf{开发者工具：}Git, Vim, VS Code
\end{itemize}

\end{document}
